% Generated by `taskgen course examples` from lattices.babai_compare (seed 6).
% Do not edit: change examples.yaml and regenerate.
\begin{example}[два базиса одной решётки]\label{ex:lattice-babai-compare}
Решётка $L \subset \mathbb{R}^2$ задана базисом $v_1 = (8, 14)$, $v_2 = (-4, 2)$. Проверить, что $u_1 = (52, 64)$, $u_2 = (-20, -26)$ --- тоже базис $L$. Применить алгоритм Бабаи (округление координат) к вектору $w = (-30, 104)$ сначала в базисе $\{v_1, v_2\}$, затем в базисе $\{u_1, u_2\}$; сравнить расстояния $\|w - r\|$ и коэффициенты Адамара базисов. Какой из найденных векторов ближе к $w$?

\par\noindent\textit{Решение.} Выражаем $u_i$ через $v_1, v_2$ (правило Крамера): $u_1 = 5 v_1 - 3 v_2$, $u_2 = -2 v_1 + 1 v_2$, матрица перехода $U$ целочисленная и $\det U = 5 \cdot 1 - (-3) \cdot (-2) = -1$, т.\,е. $U \in \mathrm{GL}_2(\mathbb{Z})$ и базисы порождают одну решётку; $\det L = 72$. Коэффициенты Адамара: $\mathcal{H}(v_1, v_2) = 0.999$ (хороший базис), $\mathcal{H}(u_1, u_2) = 0.163$ (плохой базис).\par \textbf{Базис $\{v_1, v_2\}$.} Решаем $w = t_1 v_1 + t_2 v_2$ (правило Крамера): $t_1 = \frac{89}{18} \approx 4.944$, $t_2 = \frac{313}{18} \approx 17.389$; $r = 5 v_1 + 17 v_2 = (-28, 104)$, $\|w - r\| \approx 2.000$.\par \textbf{Базис $\{u_1, u_2\}$.} $w = s_1 u_1 + s_2 u_2$: $s_1 = \frac{-715}{18} \approx -39.722$, $s_2 = \frac{-916}{9} \approx -101.778$; $r' = -40 u_1 - 102 u_2 = (-40, 92)$, $\|w - r'\| \approx 15.620$.\par Ошибка округления $w - r = \sum_i (t_i - \lfloor t_i \rceil) b_i$ с $|t_i - \lfloor t_i \rceil| \le \frac12$ ограничена величиной $\frac12(\|b_1\| + \|b_2\|)$: $\approx 10.3$ для $\{v_1, v_2\}$ и $\approx 57.6$ для длинных почти параллельных $u_1, u_2$. Вектор $r$ ближе к $w$, и это ближайший вектор решётки (сравните с $r \pm v_1$, $r \pm v_2$, $r \pm v_1 \pm v_2$); $r' = r - 1 v_1 + 1 v_2$ --- другой, далёкий вектор.\[\begin{aligned} r &= (-28, 104), & \|w - r\| &\approx 2.000;\\ r' &= (-40, 92), & \|w - r'\| &\approx 15.620.\end{aligned}\]
\end{example}
